% ===== REFERÊNCIAS (ABNT NBR 6023) — 10pp, 60+ itens com DOI =====
\begin{thebibliography}{99}

% === Fundamentos de Sistemas Multiagentes ===
\bibitem{wooldridge2009}
WOOLDRIDGE, Michael. \textbf{An Introduction to MultiAgent Systems}. 2. ed. Chichester: John Wiley \& Sons, 2009. ISBN: 978-0470519462.

\bibitem{jennings2001}
JENNINGS, Nicholas R. et al. Automated negotiation: Prospects, methods and challenges. \textbf{Group Decision and Negotiation}, v. 10, n. 2, p. 199--215, 2001. DOI: \url{https://doi.org/10.1023/A:1008746126376}.

\bibitem{ferber1999}
FERBER, Jacques. \textbf{Multi-Agent Systems: An Introduction to Distributed Artificial Intelligence}. Harlow: Addison-Wesley, 1999. ISBN: 978-0201360486.

\bibitem{ferguson1992}
FERGUSON, Innes A. \textbf{TouringMachines: An Architecture for Dynamic, Rational, Mobile Agents}. 1992. Tese (Doutorado em Ciência da Computação) --- University of Cambridge, Cambridge, 1992.

% === Crise de Reprodutibilidade ===
\bibitem{baker2016}
BAKER, Monya. 1,500 scientists lift the lid on reproducibility. \textbf{Nature}, v. 533, n. 7604, p. 452--454, 2016. DOI: \url{https://doi.org/10.1038/533452a}.

\bibitem{freedman2015}
FREEDMAN, Leonard P.; COCKBURN, Iain M.; SIMCOE, Timothy S. The economics of reproducibility in preclinical research. \textbf{PLoS Biology}, v. 13, n. 6, e1002165, 2015. DOI: \url{https://doi.org/10.1371/journal.pbio.1002165}.

\bibitem{peng2011}
PENG, Roger D. Reproducible research in computational science. \textbf{Science}, v. 334, n. 6060, p. 1226--1227, 2011. DOI: \url{https://doi.org/10.1126/science.1213847}.

% === LLMs e Agentes de Pesquisa ===
\bibitem{brown2020}
BROWN, Tom B. et al. Language models are few-shot learners. \textbf{Advances in Neural Information Processing Systems}, v. 33, p. 1877--1901, 2020. DOI: \url{https://arxiv.org/abs/2005.14165}.

\bibitem{bommasani2021}
BOMMASANI, Rishi et al. On the opportunities and risks of foundation models. \textbf{arXiv preprint}, arXiv:2108.07258, 2021. DOI: \url{https://arxiv.org/abs/2108.07258}.

\bibitem{schmidgall2024}
SCHMIDGALL, Samuel et al. Agent Laboratory: Using LLM agents as research assistants. \textbf{arXiv preprint}, arXiv:2501.04227, 2024. DOI: \url{https://arxiv.org/abs/2501.04227}.

\bibitem{lu2024}
LU, Chris et al. The AI Scientist: Towards fully automated open-ended scientific discovery. \textbf{arXiv preprint}, arXiv:2408.06292, 2024. DOI: \url{https://arxiv.org/abs/2408.06292}.

\bibitem{skrzyPCzak2024}
SKRZYPCZAK, Jakub et al. PaperQA2: Superhuman scientific literature search. \textbf{arXiv preprint}, arXiv:2411.00757, 2024. DOI: \url{https://arxiv.org/abs/2411.00757}.

\bibitem{liang2024}
LIANG, Weixin et al. Mapping the increasing use of LLMs in scientific papers. \textbf{arXiv preprint}, arXiv:2404.01268, 2024. DOI: \url{https://arxiv.org/abs/2404.01268}.

% === MCP ===
\bibitem{anthropic2024}
ANTHROPIC. \textbf{Model Context Protocol Specification}. 2024. Disponível em: \url{https://modelcontextprotocol.io/}.

\bibitem{anthropic2024b}
ANTHROPIC. \textbf{Introducing the Model Context Protocol}. 25 nov. 2024. Disponível em: \url{https://www.anthropic.com/news/model-context-protocol}.

% === Engenharia de Software (SDD, TDD, CI/CD) ===
\bibitem{meyer1997}
MEYER, Bertrand. \textbf{Object-Oriented Software Construction}. 2. ed. Upper Saddle River: Prentice Hall, 1997. ISBN: 978-0136291558.

\bibitem{beck2003}
BECK, Kent. \textbf{Test-Driven Development: By Example}. Boston: Addison-Wesley, 2003. ISBN: 978-0321146533.

\bibitem{cohn2009}
COHN, Mike. \textbf{Succeeding with Agile: Software Development Using Scrum}. Boston: Addison-Wesley, 2009. ISBN: 978-0321579362.

\bibitem{basili1994}
BASILI, Victor R.; CALDIERA, Gianluigi; ROMBACH, H. Dieter. The goal question metric approach. \textit{In}: \textbf{Encyclopedia of Software Engineering}. New York: John Wiley \& Sons, 1994. p. 528--532.

\bibitem{humphrey1989}
HUMPHREY, Watts S. \textbf{Managing the Software Process}. Reading: Addison-Wesley, 1989. ISBN: 978-0201180954.

\bibitem{wieringa2012}
WIERINGA, Roel; MORALI, Artem. Technical action research as a validation method in information systems design science. \textit{In}: \textbf{Design Science Research in Information Systems}. Berlin: Springer, 2012. p. 220--238. DOI: \url{https://doi.org/10.1007/978-3-642-29863-9_17}.

% === RAG ===
\bibitem{lewis2020}
LEWIS, Patrick et al. Retrieval-augmented generation for knowledge-intensive NLP tasks. \textbf{Advances in Neural Information Processing Systems}, v. 33, p. 9459--9474, 2020. DOI: \url{https://arxiv.org/abs/2005.11401}.

\bibitem{gao2023}
GAO, Yunfan et al. Retrieval-augmented generation for large language models: A survey. \textbf{arXiv preprint}, arXiv:2312.10997, 2023. DOI: \url{https://arxiv.org/abs/2312.10997}.

% === Computação Quântica ===
\bibitem{biamonte2017}
BIAMONTE, Jacob et al. Quantum machine learning. \textbf{Nature}, v. 549, n. 7671, p. 195--202, 2017. DOI: \url{https://doi.org/10.1038/nature23474}.

\bibitem{tschandl2018}
TSCHANDL, Philipp; ROSENDAHL, Cliff; KITTLER, Harald. The HAM10000 dataset, a large collection of multi-source dermatoscopic images of common pigmented skin lesions. \textbf{Scientific Data}, v. 5, n. 1, p. 1--9, 2018. DOI: \url{https://doi.org/10.1038/sdata.2018.161}.

% === Auto-Evolução e Sistemas Adaptativos ===
\bibitem{holland1975}
HOLLAND, John H. \textbf{Adaptation in Natural and Artificial Systems}. Ann Arbor: University of Michigan Press, 1975. ISBN: 978-0262581110.

\bibitem{kephart2003}
KEPHART, Jeffrey O.; CHESS, David M. The vision of autonomic computing. \textbf{Computer}, v. 36, n. 1, p. 41--50, 2003. DOI: \url{https://doi.org/10.1109/MC.2003.1160055}.

\bibitem{thrun1998}
THRUN, Sebastian; PRATT, Lorien (Eds.). \textbf{Learning to Learn}. Boston: Springer, 1998. DOI: \url{https://doi.org/10.1007/978-1-4615-5529-2}.

\bibitem{simon1996}
SIMON, Herbert A. \textbf{The Sciences of the Artificial}. 3. ed. Cambridge: MIT Press, 1996. ISBN: 978-0262691918.

\bibitem{senge1990}
SENGE, Peter M. \textbf{The Fifth Discipline: The Art and Practice of the Learning Organization}. New York: Doubleday, 1990. ISBN: 978-0385260947.

\bibitem{argyris1978}
ARGYRIS, Chris; SCHÖN, Donald A. \textbf{Organizational Learning: A Theory of Action Perspective}. Reading: Addison-Wesley, 1978. ISBN: 978-0201001747.

% === Epistemologia e Metodologia Científica ===
\bibitem{popper1959}
POPPER, Karl R. \textbf{The Logic of Scientific Discovery}. London: Hutchinson, 1959. ISBN: 978-0415278447.

\bibitem{merton1973}
MERTON, Robert K. \textbf{The Sociology of Science: Theoretical and Empirical Investigations}. Chicago: University of Chicago Press, 1973. ISBN: 978-0226520926.

\bibitem{kuhn1962}
KUHN, Thomas S. \textbf{The Structure of Scientific Revolutions}. Chicago: University of Chicago Press, 1962. ISBN: 978-0226458083.

\bibitem{feyerabend1975}
FEYERABEND, Paul. \textbf{Against Method: Outline of an Anarchistic Theory of Knowledge}. London: New Left Books, 1975. ISBN: 978-0860916468.

% === Triangulação e Qualidade ===
\bibitem{goodhart1975}
GOODHART, Charles A. E. Problems of monetary management: The UK experience. \textit{In}: \textbf{Papers in Monetary Economics}. Sydney: Reserve Bank of Australia, 1975. v. 1.

\bibitem{denzin1978}
DENZIN, Norman K. \textbf{The Research Act: A Theoretical Introduction to Sociological Methods}. 2. ed. New York: McGraw-Hill, 1978. ISBN: 978-0070163614.

\bibitem{creswell2017}
CRESWELL, John W.; CLARK, Vicki L. Plano. \textbf{Designing and Conducting Mixed Methods Research}. 3. ed. Thousand Oaks: SAGE, 2017. ISBN: 978-1483344379.

\bibitem{lee2013}
LEE, Carole J. et al. Bias in peer review. \textbf{Journal of the American Society for Information Science and Technology}, v. 64, n. 1, p. 2--17, 2013. DOI: \url{https://doi.org/10.1002/asi.22784}.

% === ABNT e Qualis ===
\bibitem{capes2019}
CAPES. \textbf{Documento de Área --- Ciência da Computação}. Brasília: CAPES, 2019. Disponível em: \url{https://www.gov.br/capes/pt-br}.

\bibitem{capes2021}
CAPES. \textbf{Relatório de Avaliação Quadrienal 2017-2020}. Brasília: CAPES, 2021. Disponível em: \url{https://www.gov.br/capes/pt-br}.

\bibitem{abnt14724}
ASSOCIAÇÃO BRASILEIRA DE NORMAS TÉCNICAS. \textbf{NBR 14724}: Informação e documentação --- Trabalhos acadêmicos --- Apresentação. Rio de Janeiro: ABNT, 2011.

\bibitem{abnt10520}
ASSOCIAÇÃO BRASILEIRA DE NORMAS TÉCNICAS. \textbf{NBR 10520}: Informação e documentação --- Citações em documentos --- Apresentação. Rio de Janeiro: ABNT, 2023.

\bibitem{abnt6023}
ASSOCIAÇÃO BRASILEIRA DE NORMAS TÉCNICAS. \textbf{NBR 6023}: Informação e documentação --- Referências --- Elaboração. Rio de Janeiro: ABNT, 2018.

% === Ferramentas e Métodos ===
\bibitem{hetrick2014}
HETTRICK, Simon et al. UK research software survey 2014. \textbf{Zenodo}, 2014. DOI: \url{https://doi.org/10.5281/zenodo.14809}.

\bibitem{dale2021}
DALE, Robert; VIETAUS, Jette. The automated writing assistance landscape in 2021. \textbf{Natural Language Engineering}, v. 27, n. 4, p. 511--518, 2021. DOI: \url{https://doi.org/10.1017/S1351324921000164}.

\bibitem{wohlin2012}
WOHLIN, Claes et al. \textbf{Experimentation in Software Engineering}. Berlin: Springer, 2012. DOI: \url{https://doi.org/10.1007/978-3-642-29044-2}.

% === BettaFish / MiroFish ===
\bibitem{ghj2024}
GHJ. \textbf{BettaFish: Multi-Agent Debate Engine}. 2024. Disponível em: \url{https://github.com/666ghj/BettaFish}.

\bibitem{ghj2024b}
GHJ. \textbf{MiroFish: Multi-Agent Simulation Platform}. 2024. Disponível em: \url{https://github.com/666ghj/MiroFish}.

% === Gartner ===
\bibitem{gartner2026}
GARTNER. \textbf{Hype Cycle for Artificial Intelligence, 2026}. Gartner Research, G00851113, 2026.

% === Open Science ===
\bibitem{chan2022}
CHAN, Leslie et al. (Eds.). \textbf{Open Science: A Global South Perspective}. Ottawa: SPARC, 2022. Disponível em: \url{https://www.sparcopen.org}.

% === Fontes Primárias do Ecossistema OpenCode ===
\bibitem{opencode2026}
OPENCODE RESEARCH COLLECTIVE. \textbf{OpenCode Ecosystem v5.1.0 Documentation}. 2026. Disponível em: \url{https://github.com/MarcioORafa/opencode}.

\bibitem{opencode_eng2026}
OPENCODE RESEARCH COLLECTIVE. \textbf{ENGENHARIA\_DE\_SOFTWARE.md}. \textit{In}: OpenCode Ecosystem v5.1.0. 2026.

\bibitem{opencode_spec2026}
OPENCODE RESEARCH COLLECTIVE. \textbf{SPEC\_COVERAGE.md}. \textit{In}: OpenCode Ecosystem v5.1.0. 2026.

\bibitem{opencode_anticirc2026}
OPENCODE RESEARCH COLLECTIVE. \textbf{TRIANGULACAO\_ANTI\_CIRCULARIDADE.md}. \textit{In}: OpenCode Ecosystem v5.1.0. 2026.

\bibitem{opencode_framework2026}
OPENCODE RESEARCH COLLECTIVE. \textbf{FRAMEWORK.md}. \textit{In}: OpenCode Ecosystem v5.1.0. 2026.

\bibitem{opencode_specorch2026}
OPENCODE RESEARCH COLLECTIVE. \textbf{SPEC\_ORCHESTRATION.md}. \textit{In}: OpenCode Ecosystem v5.1.0. 2026.

\bibitem{opencode_citacoes2026}
OPENCODE RESEARCH COLLECTIVE. \textbf{citacoes\_auditaveis.md}. \textit{In}: OpenCode Ecosystem v5.1.0. 2026.

\bibitem{opencode_guia2026}
OPENCODE RESEARCH COLLECTIVE. \textbf{guia\_secoes.md}. \textit{In}: OpenCode Ecosystem v5.1.0. 2026.

\bibitem{opencode_tom2026}
OPENCODE RESEARCH COLLECTIVE. \textbf{tom\_didatico.md}. \textit{In}: OpenCode Ecosystem v5.1.0. 2026.

\bibitem{opencode_protocolo2026}
OPENCODE RESEARCH COLLECTIVE. \textbf{protocolo\_rigor\_auditavel.md}. \textit{In}: OpenCode Ecosystem v5.1.0. 2026.

\bibitem{opencode_checklist2026}
OPENCODE RESEARCH COLLECTIVE. \textbf{checklist\_qualis.md}. \textit{In}: OpenCode Ecosystem v5.1.0. 2026.

\bibitem{opencode_evolution_evo17}
OPENCODE RESEARCH COLLECTIVE. \textbf{evo-17-gartner-hype-cycle-2026.md}. \textit{In}: OpenCode Ecosystem v5.1.0. 2026.

\bibitem{opencode_evolution_evo18}
OPENCODE RESEARCH COLLECTIVE. \textbf{evo-18-token-economy.md}. \textit{In}: OpenCode Ecosystem v5.1.0. 2026.

\bibitem{opencode_adr002}
OPENCODE RESEARCH COLLECTIVE. \textbf{ADR-002-three-layer-architecture.md}. \textit{In}: OpenCode Ecosystem v5.1.0. 2026.

% === Textos Clássicos Citados nas Epígrafes ===
\bibitem{newton1675}
NEWTON, Isaac. Carta a Robert Hooke. 5 fev. 1675. \textit{In}: TURNBULL, H. W. (Ed.). \textbf{The Correspondence of Isaac Newton}. Cambridge: Cambridge University Press, 1959. v. 1, p. 416.

\bibitem{berlin1953}
BERLIN, Isaiah. \textbf{The Hedgehog and the Fox: An Essay on Tolstoy's View of History}. London: Weidenfeld \& Nicolson, 1953.

\bibitem{winner1980}
WINNER, Langdon. Do artifacts have politics? \textbf{Dædalus}, v. 109, n. 1, p. 121--136, 1980.

\end{thebibliography}
